\clearpage
\VerdierPhysicalPageAnchor{0046}%
\VerdierSemanticAnchor{ch1-1-4-3-canonical}
这个作用称为\emph{典范作用}。我们要研究那些 $H$-函子

\VerdierSemanticAnchor{ch1-1-4-3-eq3}
\begin{equation}
  (F,m)\colon(\mathcal C,H,T_f,c_f)
  \longrightarrow(\mathcal C',H,T'_f,c'_f)
  \tag{1.4.3.3}
\end{equation}
它们在函子 $F$ 上诱导群 $\mathbb Z/2\mathbb Z$ 的典范作用。

\VerdierSemanticAnchor{ch1-1-4-3-extension}
设 $\dot\alpha\in H^2(G,\mathbb Z/2\mathbb Z)$ 为描述扩张

\VerdierSemanticAnchor{ch1-1-4-3-eq4}
\begin{equation}
  1\longrightarrow\mathbb Z/2\mathbb Z
  \xrightarrow{n}H\xrightarrow{f}G\longrightarrow e
  \tag{1.4.3.4}
\end{equation}
的元素。

\VerdierSemanticAnchor{ch1-1-4-3-cocycle}
设 $\alpha_s\in Z^2(G,\mathbb Z/2\mathbb Z)$ 为类 $\dot\alpha$ 中的一个
正规化非齐次余循环，它由 $f$ 的一个正规化截面 $s$ 定义。于是

\VerdierSemanticAnchor{ch1-1-4-3-eq5}
\begin{equation}
  \begin{aligned}
    \alpha_s(\sigma_1,\sigma_2)
      &=s(\sigma_2)s(\sigma_1\sigma_2)^{-1}s(\sigma_1),
      & \sigma_1,\sigma_2&\in G,\\
    s(e)&=H\text{ 的单位元},\\
    \alpha_s(\sigma_1,e)&=\alpha_s(e,\sigma_1)=1,\\
    \delta\alpha_s(\sigma_1,\sigma_2,\sigma_3)
      &=\alpha_s(\sigma_2,\sigma_3)
        \alpha_s(\sigma_1\sigma_2,\sigma_3)^{-1}\\
      &\quad{}\alpha_s(\sigma_1,\sigma_2\sigma_3)
        \alpha_s(\sigma_1,\sigma_2)^{-1}=1,
      & \sigma_1,\sigma_2,\sigma_3&\in G.
  \end{aligned}
  \tag{1.4.3.5}
\end{equation}

\VerdierSemanticAnchor{ch1-1-4-3-bijection}
截面 $s\colon G\to H$ 定义集合的一个同构
\[
  H\xrightarrow{\sim}\mathbb Z/2\mathbb Z\times G.
\]
通过这个同构把 $H$ 的群律搬到 $\mathbb Z/2\mathbb Z\times G$ 上，得到群律

\VerdierSemanticAnchor{ch1-1-4-3-eq6}
\begin{equation}
  (\varepsilon_1,\sigma_1)(\varepsilon_2,\sigma_2)
  =\bigl(\varepsilon_1\varepsilon_2
    \alpha_s(\sigma_1,\sigma_2),\sigma_1\sigma_2\bigr).
  \tag{1.4.3.6}
\end{equation}

\VerdierSemanticAnchor{ch1-1-4-3-identification}
现在把 $H$ 与 (1.4.3.6) 所描述的群等同。对任意
$(\varepsilon_1,\sigma_1),(\varepsilon_2,\sigma_2)\in H$，有

\VerdierSemanticAnchor{ch1-1-4-3-eq7}
\begin{equation}
  \begin{aligned}
    T_f(\varepsilon_1,\sigma_1)&=T(\sigma_1),\\
    T'_f(\varepsilon_1,\sigma_1)&=T'(\sigma_1),\\
    c_f\bigl((\varepsilon_1,\sigma_1),(\varepsilon_2,\sigma_2)\bigr)
      &=c(\sigma_1,\sigma_2),\\
    c'_f\bigl((\varepsilon_1,\sigma_1),(\varepsilon_2,\sigma_2)\bigr)
      &=c'(\sigma_1,\sigma_2).
  \end{aligned}
  \tag{1.4.3.7}
\end{equation}
